\chapter{Introduzione}
\label{chap:introduzione}
%
L'innovazione tecnologica nel campo dei sistemi distribuiti e della robotica collettiva ha reso sempre più evidente la necessità di un paradigma in grado di astrarre la complessità derivante dalla coordinazione di dispositivi eterogenei. In questo contesto si inserisce l'\textbf{Aggregate Computing}, un approccio che modella una rete di dispositivi interagenti come un'unica entità computazionale, consentendo di definire comportamenti collettivi attraverso costrutti dichiarativi e robusti. Questo paradigma, fondato sul \textit{field calculus}, offre gli strumenti per progettare sistemi scalabili e auto-adattivi, in cui il comportamento globale emerge da interazioni locali, senza la necessità di un controllo centralizzato.

Il \textbf{Project Emerge}, contesto applicativo del presente lavoro, rappresenta un'applicazione concreta di queste idee: un dimostratore pratico e open-source per il coordinamento di squadre di robot mobili, presentato in occasione della Notte Europea dei Ricercatori. Il progetto, basato sulla libreria \textit{MacroSwarm} e sul framework \textit{ScaFi}, mira a colmare il divario tra la teoria dell'Aggregate Computing e la sua applicazione su hardware reale, dimostrando capacità come la formazione spaziale autonoma, il \textit{self-healing} e l'indipendenza dalla topologia della rete. Tuttavia, l'architettura iniziale del sistema presentava un limite strutturale: la dipendenza rigida dal protocollo \ac{MQTT} come unico mezzo di comunicazione tra i moduli. Questo vincolo rendeva impossibile l'integrazione di robot o client che non supportassero nativamente \ac{MQTT}, limitando la flessibilità e l'estendibilità del sistema in scenari reali, dove l'eterogeneità dei dispositivi e dei protocolli è la norma.

L'obiettivo di questo lavoro di tesi è stato quindi quello di superare tale limitazione attraverso la progettazione e l'implementazione di un \textbf{gateway IoT estendibile}, in grado di disaccoppiare il protocollo interno del sistema da quelli utilizzati per comunicare con l'esterno. Il gateway, denominato \textit{domain gateway}, funge da ponte tra il broker \ac{MQTT} interno e i client esterni, traducendo i messaggi tra protocolli eterogenei (come \ac{HTTP} \ac{REST}, WebSocket e \ac{CoAP}) e garantendo che il nucleo del sistema rimanga invariato. Questo approccio non solo consente l'integrazione di dispositivi e client non compatibili con \ac{MQTT}, ma introduce anche una separazione netta tra il mondo interno e quello esterno, aprendo la strada a future evoluzioni, come la sostituzione del protocollo interno senza impattare sui client già esistenti.

Il documento è organizzato come segue. Nel \cref{chap:background} vengono presentati i fondamenti teorici dell'Aggregate Computing e l'architettura del \textbf{Project Emerge}, evidenziandone i limiti che hanno motivato questo lavoro. Il \cref{chap:analisi} formalizza il problema, definendo obiettivi, vincoli, requisiti funzionali e non funzionali, e introducendo il modello di dominio che ha guidato le scelte progettuali. Il \cref{chap:progettazione} descrive l'architettura del gateway, sia a livello di sistema che di dettaglio, illustrando le soluzioni adottate per garantire estendibilità, responsività e manutenibilità. Il \cref{chap:implementazione} approfondisce le decisioni implementative, dallo stack tecnologico all'organizzazione del codice, fino ai dettagli di ciascuna connessione protocollare. Infine, il \cref{chap:valutazione} presenta la strategia di testing automatico adottata per validare il sistema, mentre il \cref{chap:conclusione} riassume i risultati raggiunti e traccia possibili sviluppi futuri.
%